\documentclass[11pt,a4paper]{book}

\usepackage[margin=25mm]{geometry}
\usepackage{lmodern}
\usepackage{enumitem}
\usepackage{xcolor}
\usepackage{hyperref}
\hypersetup{colorlinks=true,linkcolor=blue!50!black,urlcolor=blue!50!black}

\newcommand{\abstractentry}[4]{%
  \chapter*{#1}\addcontentsline{toc}{chapter}{#1}%
  \textbf{#2}\hfill\textit{#3}\par\smallskip
  \textit{Keywords:} #4\par\medskip
}

\title{Book of Abstracts}
\author{Conference or Workshop Name}
\date{\today}

\begin{document}
\frontmatter
\maketitle
\chapter*{Welcome}
This starter is organised for a conference book of abstracts. Replace the
sample entries with your programme data, or generate this section from a CSV
file with a small script.
\tableofcontents

\mainmatter
\abstractentry{A reproducible workflow}{Alex Morgan}{Department of Computing}
  {open science, workflows, documentation}
Abstracts can be kept short and consistent while still allowing mathematical
notation such as $E=mc^2$ and references to the event programme.

\abstractentry{Designing inclusive meetings}{Jordan Lee and Sam Patel}{Institute for Learning}
  {accessibility, collaboration, teaching}
Use one command per contribution so that names, affiliations, and keywords are
easy to update when the programme changes.

\chapter*{Author index}
\addcontentsline{toc}{chapter}{Author index}
\begin{description}[leftmargin=!,labelwidth=35mm]
  \item[Morgan, Alex] A reproducible workflow
  \item[Lee, Jordan] Designing inclusive meetings
  \item[Patel, Sam] Designing inclusive meetings
\end{description}

\chapter*{Keyword index}
\addcontentsline{toc}{chapter}{Keyword index}
\begin{description}[leftmargin=!,labelwidth=35mm]
  \item[accessibility] Designing inclusive meetings
  \item[open science] A reproducible workflow
  \item[workflows] A reproducible workflow
\end{description}

\backmatter
\chapter*{Organiser notes}
The source is deliberately plain: add a bibliography, logos, or a generated
data import only when your event has cleared the relevant rights.
\end{document}
